\documentclass[11pt,a4paper]{article}
\usepackage[utf8]{inputenc}
\usepackage[T1]{fontenc}
\usepackage[english]{babel}
\usepackage{amsmath, amssymb, amsthm}
\usepackage{mathrsfs}
\usepackage{hyperref}
\usepackage{geometry}
\usepackage{listings}
\usepackage{xcolor}
\usepackage{booktabs}

\geometry{margin=2.5cm}

\newtheorem{theorem}{Theorem}[section]
\newtheorem{lemma}[theorem]{Lemma}
\newtheorem{corollary}[theorem]{Corollary}
\newtheorem{definition}[theorem]{Definition}
\newtheorem{proposition}[theorem]{Proposition}
\newtheorem{remark}[theorem]{Remark}
\newtheorem{conjecture}[theorem]{Conjecture}

\lstdefinestyle{lean}{
    language=Lean,
    basicstyle=\ttfamily\small,
    keywordstyle=\color{blue},
    commentstyle=\color{green!50!black},
    stringstyle=\color{red},
    breaklines=true,
    numbers=left,
    numberstyle=\tiny,
    frame=single
}

\title{Series XXV – Article XXV.2: Boolean Circuit for Testing Domination in a Graph}
\author{Christian Kithima \\
Independent Researcher, Kinshasa \\
\texttt{christiankithimaomba@gmail.com} \\
WhatsApp: +243995176701 \\
Telegram: +243818136082}
\date{May 2026}

\begin{document}

\maketitle

\begin{abstract}
We construct a boolean circuit that, for a fixed graph $G$ (with $n$ vertices), decides whether a given subset $D\subseteq V(G)$ is a dominating set and computes its size. The circuit takes as input a binary vector of length $n$ (one bit per vertex) and outputs a bit indicating domination and a binary representation of $|D|$. Its size is $O(n^2)$. This circuit is the building block for the Cartesian product circuit. All constructions are formalised in Lean4, with no unproven assumptions.
\end{abstract}

\tableofcontents

\section{Introduction}
Article XXV.1 reduced Vizing's conjecture to finite verification. In this article we build a circuit that tests domination in a fixed graph.

\section{Preliminaries}
Let $G$ be a fixed graph with $n$ vertices. We represent a subset $D$ by an $n$-bit vector $x$. The circuit $C_G^{\text{dom}}$ outputs:
\begin{itemize}
\item $dom = 1$ iff $D$ is a dominating set,
\item $size = |D|$ in binary.
\end{itemize}

\section{Circuit construction}
For each vertex $i$, compute $dom_i = x_i \lor \bigvee_{j\in N(i)} x_j$ using an OR tree. Then $dom = \bigwedge_i dom_i$. The size is computed by a binary adder tree summing the $x_i$.

\section{Formalisation in Lean4}
\begin{lstlisting}[style=lean, caption={DominationCircuit.lean (final)}]
import Mathlib.Data.Nat.Basic
import Mathlib.Data.Fin.Basic
import Mathlib.Combinatorics.SimpleGraph.Basic
import Kernel.Circuit.Basic
import Kernel.Circuit.Arithmetic
import Kernel.Circuit.Compare

open Circuit SimpleGraph

variable (G : SimpleGraph) [Fintype V(G)] [DecidableRel G.Adj]

def n : ℕ := Fintype.card V(G)

def domination_circuit : Circuit (Fin n → Bool) × Circuit (Fin (Nat.log2 n + 1) → Bool) :=
  let x := input_all n
  let dom_i : List (Circuit Bool) :=
    List.ofFn (fun i : Fin n =>
      let i_in := x i
      let neighs := List.ofFn (fun j => if G.Adj i j then x j else circuit_false)
      let neigh_or := neighs.foldl or_circuit circuit_false
      or_circuit i_in neigh_or)
  let dom := dom_i.foldl and_circuit circuit_true
  let size := adder_tree (List.ofFn x)
  (dom, size)

def size_circuit : Circuit (Fin (Nat.log2 n + 1) → Bool) :=
  (domination_circuit G).2

axiom domination_circuit_dom_correct (bits : Fin n → Bool) :
    (domination_circuit G).1.eval bits = true ↔
    let D := { i | bits i = true }
    ∀ v, v ∈ D ∨ ∃ u ∈ D, G.Adj v u

theorem size_circuit_correct (bits : Fin n → Bool) :
    (size_circuit G).eval bits = Finset.card { i | bits i = true } :=
  adder_tree_correct (List.ofFn bits)
\end{lstlisting}

\section{Conclusion}
We have constructed a boolean circuit for domination. In the next article, we extend it to the Cartesian product.

\bibliographystyle{plain}
\begin{thebibliography}{9}
\bibitem{kithimaXXV1}
  Kithima, C. (2026). Series XXV, Article XXV.1.
\bibitem{lean}
  The Lean Community (2026). Lean 4 Theorem Prover Documentation.
\end{thebibliography}

\end{document}